\begin{thebibliography}{99}
\bibitem{aco1}
M.~Dorigo, V.~Maniezzo y A.~Colorni. «Ant System: Optimization by a colony of
        cooperating agents». IEEE Trans. Systems, Man, and Cybernetics, Part B,
        26(1):29--41, 1996.
        \url{https://doi.org/10.1109/3477.484436}\\
        Respalda: regla de transición probabilística, evaporación y depósito del
        Ant System.

\bibitem{aco2}
M.~Dorigo y L.~M.~Gambardella. «Ant Colony System: A cooperative learning approach to
        the traveling salesman problem». IEEE Trans. Evolutionary Computation,
        1(1):53--66, 1997.
        \url{https://doi.org/10.1109/4235.585892}\\
        Respalda: variante ACS, actualización local y parámetro $q_0$.

\bibitem{aco3}
M.~Dorigo y T.~Stützle. Ant Colony Optimization. MIT Press, 2004.
        \url{https://mitpress.mit.edu/9780262042192/ant-colony-optimization/}\\
        Respalda: visión general de ACO, estigmergia, MMAS, valores de parámetros
        recomendados para el TSP y combinación con búsqueda local.

\bibitem{aco4}
P.~Larrañaga, C.~M.~H.~Kuijpers, R.~H.~Murga, I.~Inza y S.~Dizdarevic. «Genetic
        algorithms for the travelling salesman problem: A review of representations and
        operators». Artificial Intelligence Review, 13(2):129--170, 1999.
        \url{https://doi.org/10.1023/A:1006529012972}\\
        Respalda: representaciones de rutas (camino, adyacencia, ordinal) y su
        efecto en la búsqueda.

\bibitem{aco5}
C.~A.~Coello Coello. «Theoretical and numerical constraint-handling techniques used
        with evolutionary algorithms: a survey of the state of the art».
        Computer Methods in Applied Mechanics and Engineering, 191(11--12):1245--1287,
        2002. \url{https://doi.org/10.1016/S0045-7825(01)00323-1}\\
        Respalda: clasificación de técnicas para manejar restricciones (penalización,
        reparación, operadores que preservan factibilidad).

\bibitem{aco6}
T.~H.~Cormen, C.~E.~Leiserson, R.~L.~Rivest y C.~Stein. Introduction to
        Algorithms, 3.ª ed., MIT Press, 2009, sección 22.1.
        ISBN 978-0-262-03384-8.\\
        Respalda: definiciones de grafo y representaciones por matriz y por listas
        de adyacencia, con sus costos.

\bibitem{aco7}
R.~M.~Karp. «Reducibility among combinatorial problems». En Complexity of
        Computer Computations, pp.~85--103, Plenum Press, 1972.
        \url{https://doi.org/10.1007/978-1-4684-2001-2_9}\\
        Respalda: el ciclo hamiltoniano es NP-completo.

\bibitem{aco8}
M.~Haythorpe. «FHCP Challenge Set: The first set of structurally difficult instances
        of the Hamiltonian cycle problem». Bulletin of the ICA, 83:98--107, 2018.
        \url{https://arxiv.org/abs/1902.10352}\\
        Respalda: 1001 instancias hamiltonianas cuya dificultad proviene de la
        estructura, no solo del tamaño.

\bibitem{aco9}
Flinders Hamiltonian Cycle Project. «FHCP Challenge Set» (sitio oficial).
        \url{https://sites.flinders.edu.au/flinders-hamiltonian-cycle-project/fhcp-challenge-set/}\\
        Respalda: rango de 66 a 9528 vértices, formato HCP y TOUR, resultado del
        concurso (985 de 1001 grafos resueltos por el equipo ganador).

\bibitem{aco10}
L.~Pósa. «Hamiltonian circuits in random graphs». Discrete Mathematics,
        14(4):359--364, 1976. \url{https://doi.org/10.1016/0012-365X(76)90068-6}\\
        Respalda: técnica de rotación para extender caminos hacia ciclos hamiltonianos.

\bibitem{aco11}
J.~C.~Bean. «Genetic algorithms and random keys for sequencing and optimization».
        ORSA Journal on Computing, 6(2):154--160, 1994.
        \url{https://doi.org/10.1287/ijoc.6.2.154}\\
        Respalda: representación por llaves aleatorias.

\bibitem{aco12}
I.~A.~Wagner y A.~M.~Bruckstein. «Hamiltonian(t) -- An ant-inspired heuristic for
        recognizing Hamiltonian graphs». Proc. Congress on Evolutionary Computation
        (CEC 99), 1999. \url{https://doi.org/10.1109/CEC.1999.782656}\\
        Respalda: antecedente de hormigas aplicadas específicamente al ciclo
        hamiltoniano.

\bibitem{aco13}
T.~Stützle y H.~H.~Hoos. «MAX--MIN Ant System». Future Generation Computer
        Systems, 16(8):889--914, 2000.
        \url{https://doi.org/10.1016/S0167-739X(00)00043-1}\\
        Respalda: depósito solo de la mejor hormiga y límites de feromona.

\bibitem{abc1}
Flinders Hamiltonian Cycle Project, ``Flinders Hamiltonian Cycle Project'', s.~f., \url{https://sites.flinders.edu.au/flinders-hamiltonian-cycle-project/}. Define el problema hamiltoniano y distingue su objetivo del TSP. El equipo del proyecto, ``FHCP Challenge Set'', s.~f., \url{https://sites.flinders.edu.au/flinders-hamiltonian-cycle-project/fhcp-challenge-set/}, describe el conjunto FHCP y su formato. Se citan como fuente institucional de contexto, no como evidencia experimental de este ejemplo.

\bibitem{abc2}
Michael Haythorpe, ``FHCP Challenge Set: The First Set of Structurally Difficult Instances of the Hamiltonian Cycle Problem'', 2018, \url{https://arxiv.org/abs/1902.10352}. Respalda la identidad y finalidad del conjunto FHCP. La fecha de publicación del artículo de revista es 2018; la versión arXiv enlazada es posterior.

\bibitem{abc3}
David L. Poole y Alan K. Mackworth, Artificial Intelligence: Foundations of Computational Agents, 3.\textsuperscript{a} ed., Cambridge University Press, 2023, sección 3.6, \url{https://www.artint.info/3e/html/ArtInt3e.Ch3.S6.html}. Respalda la explicación de búsqueda informada y de función heurística como guía estimativa.

\bibitem{abc4}
John H. Holland, Adaptation in Natural and Artificial Systems, 1975; reimpresión MIT Press, 1992, \url{https://doi.org/10.7551/MITPRESS/1090.001.0001}. Fundamento histórico del algoritmo genético y de la adaptación mediante recombinación y selección.

\bibitem{abc5}
Xin-She Yang, ``Metaheuristic Optimization'', Scholarpedia, 2011, \url{https://www.scholarpedia.org/article/Metaheuristic_Optimization}. Respalda el esquema general de GA: codificación, fitness, población, selección, cruce, mutación y reemplazo.

\bibitem{abc6}
David E. Goldberg y Robert Lingle Jr., ``Alleles, Loci, and the Traveling Salesman Problem'', Proceedings of the First International Conference on Genetic Algorithms, 1985, \url{https://gpbib.pmacs.upenn.edu/icga/icga1985-ocr.pdf}. Respalda el uso de permutaciones y cruces especializados, como PMX, para recorridos.

\bibitem{abc7}
Derviş Karaboğa, ``An Idea Based on Honey Bee Swarm for Numerical Optimization'', informe técnico TR06, Universidad de Erciyes, 2005, \href{https://citeseerx.ist.psu.edu/document?doi=cf20e34a1402a115523910d2a4243929f6704db1&repid=rep1&type=pdf}{enlace al informe}. Fuente primaria de las fases de ABC, roles de abejas, fuentes, selección y parámetro \texttt{limit}.

\bibitem{abc8}
Derviş Karaboğa, ``Artificial Bee Colony Algorithm'', Scholarpedia 5(3):6915, 2010, \href{https://www.scholarpedia.org/article/Artificial_bee_colony_algorithm}{enlace al artículo}. Expone la versión de ABC con selección probabilística, mejora voraz, abandono y controles principales.

\bibitem{abc9}
Indadul Khan y Manas Kumar Maiti, ``A Swap Sequence Based Artificial Bee Colony Algorithm for Traveling Salesman Problem'', Swarm and Evolutionary Computation 44, 428--438, 2019, \url{https://doi.org/10.1016/j.swevo.2018.05.006}. Respalda que las variantes discretas para recorridos requieren operadores de permutación; no se atribuye a este artículo la adaptación específica propuesta aquí.

\bibitem{de1}
Richard M. Karp, ``Reducibility Among Combinatorial Problems'', en Miller y Thatcher (eds.), Complexity of Computer Computations, Plenum, 85--103, 1972, \url{https://people.irisa.fr/Nicolas.Markey/biblio/ccc1972-Kar}. Marco clásico de NP-completitud. Trabajo antiguo y vigente; se verificó la referencia, y queda pendiente confirmar el pasaje específico en el texto.

\bibitem{de2}
Michael Haythorpe, ``FHCP Challenge Set: The First Set of Structurally Difficult Instances of the Hamiltonian Cycle Problem'', Bulletin of the ICA 83, 98--107, 2018, \url{https://arxiv.org/abs/1902.10352}. Fuente primaria sobre la composición y dificultad del conjunto, y sobre la adaptación de algoritmos del TSP.

\bibitem{de3}
Flinders Hamiltonian Cycle Project, ``FHCP Challenge Set'', consultado en noviembre de 2026, \url{https://sites.flinders.edu.au/flinders-hamiltonian-cycle-project/fhcp-challenge-set/}. Página oficial: tamaños, existencia de ciclo en todas las instancias y resultados del concurso.

\bibitem{de4}
Joeri Sleegers y Daan van den Berg, ``Looking for the Hardest Hamiltonian Cycle Problem Instances'', ECTA 2020, 40--48, 2020, \url{https://doi.org/10.5220/0010066900400048}. Respalda la relación entre densidad y facilidad. Artículo de congreso.

\bibitem{de5}
Thomas Bartz-Beielstein et al., ``Benchmarking in Optimization: Best Practice and Open Issues'', arXiv:2007.03488v2, 2020, \url{https://arxiv.org/abs/2007.03488}. Revisión amplia de 17 autores sobre esfuerzo, éxito, fallos, PAR-$k$, resumen de ejecuciones y semillas. Es un preprint sin arbitraje confirmado.

\bibitem{de6}
Nikolaus Hansen, Anne Auger, Dimo Brockhoff, Dejan Tušar y Tea Tušar, ``COCO: Performance Assessment'', arXiv:1605.03560, 2016, \url{https://arxiv.org/abs/1605.03560}. Respalda las evaluaciones como medida, el objetivo fijo, el tratamiento de fallos y el aRT. Su contexto es la optimización continua, de modo que aplicarlo aquí es una transferencia.

\bibitem{de7}
A. E. Eiben y J. E. Smith, Introduction to Evolutionary Computing, 2.\textsuperscript{a} ed., Springer, 2015, \url{https://doi.org/10.1007/978-3-662-44874-8}. Respalda eficiencia, efectividad y la advertencia sobre generaciones. Solo se revisaron \href{https://cin.ufpe.br/~aluizioa/CE/CE-07-trabalhando-com-AEs-ch09-2014.pdf}{diapositivas del capítulo}; el libro queda pendiente.

\bibitem{de8}
Lawrence D. Brown, T. Tony Cai y Anirban DasGupta, ``Interval Estimation for a Binomial Proportion'', Statistical Science 16(2), 101--133, 2001, \url{https://doi.org/10.1214/ss/1009213286}. Respalda la incertidumbre de una tasa de éxito. Se revisó el resumen.

\bibitem{de9}
Peter Christen, David J. Hand y Nishadi Kirielle, ``A Review of the F-Measure: Its History, Properties, Criticism, and Alternatives'', ACM Computing Surveys 56(3), artículo 73, 2023, \url{https://doi.org/10.1145/3606367}. Respalda definiciones y críticas de precisión, recall y F1. Trata la clasificación; no se le atribuye la aplicación al ciclo hamiltoniano hecha aquí.

\bibitem{de10}
Rainer Storn y Kenneth Price, ``Differential Evolution: A Simple and Efficient Heuristic for Global Optimization over Continuous Spaces'', Journal of Global Optimization 11, 341--359, 1997, \href{https://www.cp.eng.chula.ac.th/~piak/teaching/ec/ec2012/storn_price_de.pdf}{copia consultada}. Fuente primaria del algoritmo, la notación y los parámetros. El DOI queda pendiente de confirmar.

\bibitem{de11}
Swagatam Das, Sankha Subhra Mullick y P. N. Suganthan, ``Recent Advances in Differential Evolution: An Updated Survey'', Swarm and Evolutionary Computation 27, 1--30, 2016, \url{https://doi.org/10.1016/j.swevo.2016.01.004}. Respalda el desarrollo de DE y la adaptación de parámetros. Se revisó el resumen.

\bibitem{de12}
Karin Zielinski, Petra Weitkemper, Rainer Laur y Karl-Dirk Kammeyer, ``Examination of Stopping Criteria for Differential Evolution based on a Power Allocation Problem'', 10th International Conference on Optimization of Electrical and Electronic Equipment, 2006, \url{https://www.ant.uni-bremen.de/sixcms/media.php/102/4605/ICO3E_2006_zielinski.pdf}. Respalda la clasificación de criterios de paro. Un congreso y un solo problema.

\bibitem{de13}
James C. Bean, ``Genetic Algorithms and Random Keys for Sequencing and Optimization'', ORSA Journal on Computing 6(2), 154--160, 1994, \url{https://doi.org/10.1287/ijoc.6.2.154}. Origen de las claves aleatorias, propuestas para algoritmos genéticos. Se revisó el resumen.

\bibitem{de14}
Pavel Krömer, Vojtěch Uher y Václav Snášel, ``Novel Random Key Encoding Schemes for the Differential Evolution of Permutation Problems'', IEEE Transactions on Evolutionary Computation 26(1), 43--57, 2022, \url{https://doi.org/10.1109/TEVC.2021.3087802}. Respalda el uso de claves aleatorias en DE y en el TSP. Se revisó el resumen.

\bibitem{de15}
Valentino Santucci, Marco Baioletti y Alfredo Milani, ``Algebraic Differential Evolution Algorithm for the Permutation Flowshop Scheduling Problem with Total Flowtime Criterion'', IEEE Transactions on Evolutionary Computation 20(5), 682--694, 2016, \url{https://doi.org/10.1109/TEVC.2015.2507785}. Respalda la crítica a las claves aleatorias y la DE algebraica, en otro problema.

\bibitem{de16}
Godfrey C. Onwubolu y Donald Davendra (eds.), Differential Evolution: A Handbook for Global Permutation-Based Combinatorial Optimization, Springer, 2009, ISBN 978-3-540-92150-9, \url{https://www.wolfram.com/books/profile.cgi?id=7359}. Respalda la existencia de varios enfoques. Solo se revisaron la ficha y el índice.

\bibitem{de17}
Fatih Tasgetiren, Angela Chen, Gunes Gencyilmaz y Said Gattoufi, ``Smallest Position Value Approach'', capítulo 5 de \cite{de16}, 2009, \url{https://www.academia.edu/3714507/Smallest_Position_Value_Approach}. Respalda el uso de SPV con DE en un TSP generalizado. Se revisó un extracto; el detalle de la regla queda pendiente.
\bibitem{sa1} Leighton, T. y Rubinfeld, R. (2006). Mathematics for Computer Science: Notes for Recitation 6. MIT. Definiciones de grafos, recorridos y longitud. \url{https://web.mit.edu/neboat/Public/6.042/rec6.pdf}.
\bibitem{sa2} Diestel, R. (2025). Graph Theory, 6.ª ed. Springer, capítulo 10. Definiciones hamiltonianas; referencia conservada de la investigación de SA. \url{https://www.math.uni-hamburg.de/home/diestel/books/graph.theory/preview/Ch10.pdf}.
\bibitem{sa3} Universidad de Waterloo. The Traveling Salesman Problem. Definición del TSP. \url{https://www.math.uwaterloo.ca/tsp/problem/index.html}.
\bibitem{sa4} Kirkpatrick, S., Gelatt, C. D. y Vecchi, M. P. (1983). Optimization by Simulated Annealing. Science 220(4598), 671--680. Fundamento de SA. \url{https://doi.org/10.1126/science.220.4598.671}.
\bibitem{sa5} Hajek, B. (1988). Cooling Schedules for Optimal Annealing. Mathematics of Operations Research 13(2), 311--329. Condiciones de convergencia. \url{https://doi.org/10.1287/moor.13.2.311}.
\bibitem{sa7} Komosinski, M. Optimization: Simulated annealing. Poznan University of Technology. Material docente citado por el documento original para el pseudocódigo y el enfriamiento. \url{https://www.cs.put.poznan.pl/mkomosinski/lectures/optimization/SA-en.pdf}.
\bibitem{sa8} Metropolis, N., Rosenbluth, A. W., Rosenbluth, M. N., Teller, A. H. y Teller, E. (1953). Equation of State Calculations by Fast Computing Machines. Journal of Chemical Physics 21(6), 1087--1092. Antecedente de la aceptación probabilística. \url{https://doi.org/10.1063/1.1699114}.
\bibitem{pso1995} Kennedy, J. y Eberhart, R. (1995). Particle swarm optimization. Proceedings of ICNN'95, vol. 4, 1942--1948. Fuente original de PSO. \url{https://doi.org/10.1109/ICNN.1995.488968}.
\bibitem{pso1998} Shi, Y. y Eberhart, R. (1998). A modified particle swarm optimizer. IEEE International Conference on Evolutionary Computation, 69--73. Variante con peso de inercia. \url{https://doi.org/10.1109/ICEC.1998.699146}.

\end{thebibliography}
